\documentclass{article}
\usepackage[utf8]{inputenc}
\usepackage{hyperref}
\usepackage{listings}
\usepackage{xcolor}
\usepackage{graphicx}

\definecolor{codegreen}{rgb}{0,0.6,0}
\definecolor{codegray}{rgb}{0.5,0.5,0.5}
\definecolor{codepurple}{rgb}{0.58,0,0.82}
\definecolor{backcolour}{rgb}{0.95,0.95,0.92}

\lstdefinestyle{mystyle}{
    backgroundcolor=\color{backcolour},   
    commentstyle=\color{codegreen},
    keywordstyle=\color{magenta},
    numberstyle=\tiny\color{codegray},
    stringstyle=\color{codepurple},
    basicstyle=\ttfamily\footnotesize,
    breakatwhitespace=false,         
    breaklines=true,                 
    captionpos=b,                    
    keepspaces=true,                 
    numbers=left,                    
    numbersep=5pt,                  
    showspaces=false,                
    showstringspaces=false,
    showtabs=false,                  
    tabsize=2
}

\lstset{style=mystyle}

\title{PostgreSQL High Availability Solution}
\author{Database Engineering Team}
\date{\today}

\begin{document}

\maketitle

\tableofcontents

\section{System Overview}
\subsection{Architecture}
\begin{itemize}
\item Primary node (read/write)
\item 2+ replica nodes (hot standby)
\item Automated failover with Bash/Python
\item Health monitoring system
\end{itemize}

\subsection{Components}
\begin{tabular}{|l|l|}
\hline
\textbf{Component} & \textbf{Version} \\
\hline
PostgreSQL & 14+ \\
Streaming Replication & Native \\
Failover Manager & Bash Scripts \\
Monitoring & Custom Python \\
\hline
\end{tabular}

\section{Installation Guide}
\subsection{Requirements}
\begin{itemize}
\item 3+ Ubuntu/Debian 20.04/22.04 servers
\item SSH key-based access between nodes
\item Minimum 4GB RAM, 2 CPU cores per node
\end{itemize}

\subsection{Setup Commands}
\begin{lstlisting}[language=bash,caption={Primary Node Setup}]
sudo NODE_TYPE=primary bash setup_postgres_ha.sh
\end{lstlisting}

\begin{lstlisting}[language=bash,caption={Replica Node Setup}]
sudo NODE_TYPE=replica PRIMARY_IP=<primary_ip> bash setup_postgres_ha.sh
\end{lstlisting}

\section{Failover Procedures}
\subsection{Automated Failover}
\begin{enumerate}
\item Health checks fail (30s timeout)
\item Trigger file created
\item Replica promotion
\item Cluster reconfiguration
\end{enumerate}

\subsection{Manual Failover}
\begin{lstlisting}[language=bash,caption={Manual Promotion}]
sudo -u postgres pg_ctl promote -D /var/lib/postgresql/14/main
\end{lstlisting}

\section{Recovery Processes}
\subsection{Primary Recovery}
\begin{lstlisting}[language=bash,caption={Failed Primary Recovery}]
sudo -u postgres pg_rewind -D /var/lib/postgresql/14/main \
  --source-server="host=new_primary"
\end{lstlisting}

\section{Monitoring System}
\subsection{Metrics Collected}
\begin{itemize}
\item Node availability (up/down)
\item Replication lag (seconds)
\item Disk space usage (\%)
\item Connection latency (ms)
\end{itemize}

\subsection{Alert Thresholds}
\begin{tabular}{|l|l|l|}
\hline
\textbf{Metric} & \textbf{Warning} & \textbf{Critical} \\
\hline
Replication Lag & 10s & 30s \\
Disk Space & 80\% & 90\% \\
\hline
\end{tabular}

\section{Limitations}
\begin{itemize}
\item Maximum 5 nodes supported
\item Failover time 30-60 seconds
\item No automatic backup management
\end{itemize}

\section{Future Improvements}
\subsection{Short-term}
\begin{itemize}
\item Prometheus metrics endpoint
\item Slack notifications
\end{itemize}

\subsection{Long-term}
\begin{itemize}
\item Kubernetes operator
\item Web management UI
\end{itemize}

\section{Maintenance}
\subsection{Routine Checks}
\begin{itemize}
\item Daily replication verification
\item Monthly failover testing
\end{itemize}

\section{Troubleshooting}
\subsection{Common Issues}
\begin{lstlisting}[language=SQL,caption={High Replication Lag}]
SELECT pg_wal_replay_resume();
\end{lstlisting}

\begin{thebibliography}{9}
\bibitem{postgresql}
PostgreSQL Documentation,
\textit{https://www.postgresql.org/docs/}
\end{thebibliography}

\end{document}